% GENERATED FILE. Edit canonical lessons, then run npm run generate:books.
% canonical-source-hash: fnv1a64:d2898bfdc6aac0f5
% canonical-lessons: GE-C30-auge, GE-C30-ohr, GE-C30-mund, GE-C30-nase

\chapter{Eyes, Ears, Mouth, Nose}
\label{ch:ge-face-parts}
\hlchaptermodality{30}

\begin{chapteropening}
\textbf{By the end of this chapter:} \emph{I can name the eye, ear, mouth, and nose in German, and say which of the four traces to a confirmed ancient root and which does not.}

The chapter closes on Nase, cousin of both English nose and Latin nasus by separate descent: the reader produces all four face parts with articles and plurals, names Auge and Ohr as words with confirmed Proto-Indo-European roots, names Mund as inherited but not traced further, and completes Chapter 17's face by running Kopf, Auge, Ohr, Mund, Nase, Hand together. It reaches back to Chapter 17's Kopf and Hand and to Chapter 26's hören, rescuing the disputed 'sharp-eared' link between hören and Ohr.
\end{chapteropening}

\section[das Auge]{das Auge — back to the head, one part at a time}

\label{lesson:GE-C30-auge}



\begin{quote}
Chapter 17 gave you \textbf{der Kopf} and \textbf{die Hand}. The head has more in it than its own name — starting with what you see with.
\end{quote}

\subsection*{You'll want to know: das Auge}
\begin{quote}
\textbf{das Auge} — eye. Neuter.
\end{quote}

\begin{quote}
\textbf{Das ist das Auge.} — \textquotedblleft{}This/That is the eye.\textquotedblright{}
\end{quote}

The plural is \textbf{die Augen} — both eyes together, the form you will hear far more often than the singular.

\begin{sounds}
\textbf{au} is a diphthong said \emph{ow}, as in English \emph{how}: \textbf{Auge} = \emph{OW-geh}.
\end{sounds}

\begin{cousinweb}
\textbf{Auge} is Proto-Germanic \emph{\textbf{*augô}}, the direct cousin of English \textbf{eye} — both inherited, no loan in either direction, from Proto-Indo-European \emph{\textbf{*h\textsubscript{3}ek\textsuperscript{w}-}}. That root did not stay hidden in only the Germanic family: Latin turned it into \textbf{oculus} (giving English \textbf{ocular}, \textbf{binoculars}), and Greek into \textbf{ophthalmós} (giving English \textbf{ophthalmology}). Three branches of one family tree, and English ended up borrowing back two Latin and Greek reflexes of the very root it had already inherited natively as \emph{eye}.

Chapter 17's \emph{Kopf} was a container word that took over the everyday job. \emph{Auge} needed no such replacement — it has meant \textquotedblleft{}eye\textquotedblright{} the whole way down.
\end{cousinweb}

\subsection*{Your turn}
Say these aloud:

\begin{itemize}
\raggedright
  \item \textquotedblleft{}das Auge\textquotedblright{} — \emph{OW-geh}; \textquotedblleft{}die Augen\textquotedblright{} — both eyes
  \item \textquotedblleft{}Das ist das Auge.\textquotedblright{}
  \item the three-family root — \textquotedblleft{}Auge/eye, Latin oculus, Greek ophthalmós\textquotedblright{}
\end{itemize}

\begin{tcolorbox}[breakable,colback=teal!4,colframe=teal!35!black,title={Before you move on}]
Give \textquotedblleft{}eye\textquotedblright{} with its article, and its plural. (\textbf{Das Auge, die Augen}.) Is it native or borrowed? (\textbf{Native} — cousin of English \emph{eye}.) Name one English word borrowed from the \emph{same} PIE root by way of Latin or Greek. (\textbf{Ocular}, \textbf{binoculars}, or \textbf{ophthalmology}.)
\end{tcolorbox}
\section[das Ohr]{das Ohr — the part hören was named for}

\label{lesson:GE-C30-ohr}



\begin{quote}
You have the eye. Chapter 26 already gave you the verb for what the next part does — \textbf{hören}, \textquotedblleft{}to hear\textquotedblright{} — without ever naming the part itself.
\end{quote}

\subsection*{You'll want to know: das Ohr}
\begin{quote}
\textbf{das Ohr} — ear. Neuter.
\end{quote}

\begin{quote}
\textbf{Das ist das Ohr.} — \textquotedblleft{}This/That is the ear.\textquotedblright{}
\end{quote}

Plural: \textbf{die Ohren}.

\begin{sounds}
\textbf{Ohr} is a single long \emph{o}, said the way you'd say the letter itself in English: \textbf{Ohr} = \emph{OHR}, no glide, no second vowel.
\end{sounds}

\begin{cousinweb}
\textbf{Ohr} is Proto-Germanic \emph{\textbf{*auzô}}, the direct cousin of English \textbf{ear} — both from Proto-Indo-European \emph{\textbf{*h\textsubscript{2}ous-}}. Chapter 26 already flagged one disputed idea about \textbf{hören}: some accounts break its root into \textquotedblleft{}sharp\textquotedblright{} plus \textquotedblleft{}ear,\textquotedblright{} which would put this very word, \emph{Ohr}, hiding inside the verb for using it. That analysis was marked cited but unsettled there, and it stays unsettled here — a tidy story is not the same as a confirmed one.

What is confirmed: \emph{Ohr} and \emph{ear} are not two words that happen to sound alike. They are one Germanic word, unbroken since before English and German were different languages.
\end{cousinweb}

\subsection*{Your turn}
Say these aloud:

\begin{itemize}
\raggedright
  \item \textquotedblleft{}das Ohr\textquotedblright{} — a plain long \emph{o}; \textquotedblleft{}die Ohren\textquotedblright{}
  \item \textquotedblleft{}Das ist das Ohr.\textquotedblright{}
  \item the maybe-link — \textquotedblleft{}hören, Ohr — sharp-eared, unsettled\textquotedblright{}
\end{itemize}

\begin{tcolorbox}[breakable,colback=teal!4,colframe=teal!35!black,title={Before you move on}]
Give \textquotedblleft{}ear\textquotedblright{} with its article and plural. (\textbf{Das Ohr, die Ohren}.) Is it native or borrowed? (\textbf{Native} — cousin of English \emph{ear}.) What disputed idea from Chapter 26 would put this word inside \emph{hören}? (\textbf{The \textquotedblleft{}sharp-eared\textquotedblright{} root analysis} — cited, not settled.)
\end{tcolorbox}
\section[der Mund]{der Mund — where the family tree runs out}

\label{lesson:GE-C30-mund}



\begin{quote}
\emph{Auge} and \emph{Ohr} both trace cleanly to Proto-Indo-European. This one does not go nearly as far back with any certainty — and saying so honestly is the lesson.
\end{quote}

\subsection*{You'll want to know: der Mund}
\begin{quote}
\textbf{der Mund} — mouth. Masculine.
\end{quote}

\begin{quote}
\textbf{Das ist der Mund.} — \textquotedblleft{}This/That is the mouth.\textquotedblright{}
\end{quote}

Plural: \textbf{die Münder}.

\begin{sounds}
\textbf{Mund} takes a short, closed \emph{u}, as in English \emph{put}: \textbf{Mund} = \emph{MOONT} (final \emph{d} devoiced, as in \emph{Hand}).
\end{sounds}

\begin{cousinweb}
\textbf{Mund} is Proto-Germanic \emph{\textbf{*munþaz}}, and English \textbf{mouth} is that same inherited word — no loan in either direction, exactly like \emph{Ohr} and \emph{ear}.

Here the two words part company. \emph{Ohr} and \emph{Auge} both trace on to confirmed Proto-Indo-European roots. \emph{Munþaz} does not: proposed deeper connections exist, but no single one is agreed upon by standard references. Rather than hand you a tidy-sounding root that is not actually settled, this lesson stops at Proto-Germanic. An honest \textquotedblleft{}we don't fully know\textquotedblright{} is worth more than an invented ancestor.
\end{cousinweb}

\subsection*{Your turn}
Say these aloud:

\begin{itemize}
\raggedright
  \item \textquotedblleft{}der Mund\textquotedblright{} — \emph{MOONT}; \textquotedblleft{}die Münder\textquotedblright{}
  \item \textquotedblleft{}Das ist der Mund.\textquotedblright{}
  \item the honest stop — \textquotedblleft{}Mund and mouth, inherited — but no further-back root is agreed\textquotedblright{}
\end{itemize}

\begin{tcolorbox}[breakable,colback=teal!4,colframe=teal!35!black,title={Before you move on}]
Give \textquotedblleft{}mouth\textquotedblright{} with its article and plural. (\textbf{Der Mund, die Münder}.) Is it native or borrowed? (\textbf{Native} — cousin of English \emph{mouth}.) Does its root trace to a confirmed PIE ancestor, the way \emph{Auge} and \emph{Ohr} do? (\textbf{No} — the trail goes cold at Proto-Germanic.)
\end{tcolorbox}
\section[die Nase]{die Nase — the face, complete}

\label{lesson:GE-C30-nase}



\begin{quote}
\textbf{Auge. Ohr. Mund.} One part of the face is left, and it closes what Chapter 17 opened with \emph{Kopf} and never finished.
\end{quote}

\subsection*{You'll want to know: die Nase}
\begin{quote}
\textbf{die Nase} — nose. Feminine.
\end{quote}

\begin{quote}
\textbf{Das ist die Nase.} — \textquotedblleft{}This/That is the nose.\textquotedblright{}
\end{quote}

Plural: \textbf{die Nasen}.

\begin{sounds}
\textbf{Nase} opens on a long, open \emph{a}: \textbf{Nase} = \emph{NAH-zeh}, with the \emph{s} voiced to \emph{z} between two vowels — the mirror image of \emph{Hand}'s final \emph{d} devoicing to \emph{t}.
\end{sounds}

\begin{cousinweb}
\textbf{Nase} is Proto-Germanic \emph{\textbf{*nasō}}, the inherited cousin of English \textbf{nose} — and, further back, of Latin \textbf{nasus}, which gave English \textbf{nasal} and \textbf{nostril}. One Proto-Indo-European root, \emph{\textbf{*nas-}}, split into Germanic and into Latin independently; German \emph{Nase} and the Latin-derived word are family by \textbf{descent}, exactly like \emph{rot} and \emph{rouge} in Chapter 13, never by borrowing.

That gives this small face four parts with four different shapes of story: \textbf{Auge} and \textbf{Ohr} trace cleanly to confirmed PIE roots; \textbf{Mund} is inherited but its root before Proto-Germanic is not agreed upon; \textbf{Nase} traces to a root Latin also inherited, on its own branch. Four native words, four different depths of certainty — which is the honest way etymology actually looks, not four identical tidy stories.
\end{cousinweb}

\subsection*{Your turn}
Say these aloud:

\begin{itemize}
\raggedright
  \item \textquotedblleft{}die Nase\textquotedblright{} — \emph{NAH-zeh}; \textquotedblleft{}die Nasen\textquotedblright{}
  \item \textquotedblleft{}Das ist die Nase.\textquotedblright{}
  \item the whole face — \textquotedblleft{}Kopf, Auge, Ohr, Mund, Nase, Hand\textquotedblright{}
  \item the cousin — \textquotedblleft{}Nase, nose, nasus — one PIE root, two families\textquotedblright{}
\end{itemize}

\begin{tcolorbox}[breakable,colback=teal!4,colframe=teal!35!black,title={Before you move on}]
Give \textquotedblleft{}nose\textquotedblright{} with its article and plural. (\textbf{Die Nase, die Nasen}.) Which Latin word is its distant cousin, and what English word does that Latin word give? (\textbf{Nasus}, giving \textbf{nasal}.) Name all five face and head words from Chapters 17 and 30. (\textbf{Kopf, Auge, Ohr, Mund, Nase}.) Which of the four Chapter 30 words does not trace past Proto-Germanic? (\textbf{Mund}.)
\end{tcolorbox}
